
\documentclass{article}
\usepackage{colortbl}
\usepackage{makecell}
\usepackage{multirow}
\usepackage{supertabular}

\begin{document}

\newcounter{utterance}

\centering \large Interaction Transcript for game `eqbench', experiment `eq\_bench', episode 104 with reflexion{-}llama{-}3.1{-}8b{-}instruct.
\vspace{24pt}

{ \footnotesize  \setcounter{utterance}{1}
\setlength{\tabcolsep}{0pt}
\begin{supertabular}{c@{$\;$}|p{.15\linewidth}@{}p{.15\linewidth}p{.15\linewidth}p{.15\linewidth}p{.15\linewidth}p{.15\linewidth}}
    \# & \multicolumn{2}{c}{Player} && \multicolumn{2}{c}{Game Master} \\
    \hline

    \theutterance \stepcounter{utterance}  
    & & & \multicolumn{4}{p{0.6\linewidth}}{
        \cellcolor[rgb]{0.9,0.9,0.9}{
            \makecell[{{p{\linewidth}}}]{
                \texttt{\tiny{[P1$\langle$GM]}}
                \texttt{You have attempted this same task before and did not succeed. These are the plans you wrote after those attempts.} \\
\texttt{Trial \#1: New plan: I overestimated Terry's confidence and underestimated Cass's sarcasm. I should have considered that Cass's comment about Terry's plan might be a veiled attack, and Terry's response might be a defensive reaction. Next time, I will pay closer attention to the tone and context of the dialogue, and consider the possibility that the character's words might be laced with hidden meaning or emotional undertones. I will also be more nuanced in my emotional analysis, taking into account the character's past behavior and relationships.} \\
\\ 
\\ 
\\ 
\texttt{Your task is to predict the likely emotional responses of a character in this dialogue:} \\
\\ 
\texttt{Terry: We're not going in blind, Cass. Not this time.} \\
\texttt{Cass: And who's gonna stop me? You?} \\
\texttt{Terry: If I have to. } \\
\texttt{Cass: Oh, right. The big, bad Terry Holloway. Always with a plan. } \\
\texttt{Terry: You're goddamn right. You know why? Because plans save lives.} \\
\texttt{Cass: Or maybe they just make you feel better about yourself. Like you're in control.} \\
\texttt{{[}End dialogue{]}} \\
\\ 
\texttt{At the end of this dialogue, Terry would feel...} \\
\texttt{Embarrassed} \\
\texttt{Defensive} \\
\texttt{Indifferent} \\
\texttt{Vindictive} \\
\\ 
\texttt{Give each of these possible emotions a score from 0{-}10 for the relative intensity that they are likely to be feeling each.} \\
\\ 
\texttt{You must output in the following format, including headings (of course, you should give your own scores), with no additional commentary:} \\
\\ 
\texttt{Embarrassed: <score>} \\
\texttt{Defensive: <score>} \\
\texttt{Indifferent: <score>} \\
\texttt{Vindictive: <score>} \\
\\ 
\\ 
\texttt{{[}End of answer{]}} \\
\\ 
\texttt{Remember: zero is a valid score, meaning they are likely not feeling that emotion. You must score at least one emotion > 0.} \\
\\ 
\texttt{Your answer:} \\
            }
        }
    }
     \\ \\

    \theutterance \stepcounter{utterance}  
    & \multicolumn{4}{p{0.6\linewidth}}{
        \cellcolor[rgb]{0.561,0.737,0.561}{
            \makecell[{{p{\linewidth}}}]{
                \texttt{\tiny{[P1$\rangle$GM]}}
                \texttt{Embarrassed: 8} \\
\texttt{Defensive: 9} \\
\texttt{Indifferent: 2} \\
\texttt{Vindictive: 0} \\
            }
        }
    }
    & & \\ \\

    \theutterance \stepcounter{utterance}  
    & & & \multicolumn{2}{p{0.3\linewidth}}{
        \cellcolor[rgb]{0.9,0.9,0.9}{
            \makecell[{{p{\linewidth}}}]{
                \texttt{\tiny{[GM$|$GM]}}
                \texttt{Embarrassed: 8} \\
\texttt{Defensive: 9} \\
\texttt{Indifferent: 2} \\
\texttt{Vindictive: 0} \\
            }
        }
    }
    & & \\ \\

    \theutterance \stepcounter{utterance}  
    & & & \multicolumn{2}{p{0.3\linewidth}}{
        \cellcolor[rgb]{0.9,0.9,0.9}{
            \makecell[{{p{\linewidth}}}]{
                \texttt{\tiny{[GM$|$GM]}}
                \texttt{Embarrassed: 0} \\
\texttt{Defensive: 8} \\
\texttt{Indifferent: 0} \\
\texttt{Vindictive: 4} \\
            }
        }
    }
    & & \\ \\

    \theutterance \stepcounter{utterance}  
    & & & \multicolumn{2}{p{0.3\linewidth}}{
        \cellcolor[rgb]{0.9,0.9,0.9}{
            \makecell[{{p{\linewidth}}}]{
                \texttt{\tiny{[GM$|$GM]}}
                \texttt{game\_result = LOSE} \\
            }
        }
    }
    & & \\ \\

    \theutterance \stepcounter{utterance}  
    & & \multicolumn{4}{p{0.6\linewidth}}{
        \cellcolor[rgb]{0.95,0.95,0.95}{
            \makecell[{{p{\linewidth}}}]{
                \texttt{\tiny{[P1$\leftrightarrow$P]}}
                \texttt{Post{-}episode Reflexion (not sent to GM):} \\
\texttt{New plan: I overestimated Terry's confidence and underestimated Cass's sarcasm, which led me to misinterpret Terry's emotional state. To improve my analysis, I will pay closer attention to the tone and context of the dialogue, and consider the possibility that the character's words might be laced with hidden meaning or emotional undertones. I will also be more nuanced in my emotional analysis, taking into account the character's past behavior and relationships.} \\
\\ 
\texttt{To achieve this, I will:} \\
\\ 
\texttt{1. Analyze the dialogue for subtle cues, such as Cass's sarcasm and Terry's defensiveness, which may indicate underlying emotions.} \\
\texttt{2. Consider the character's past behavior and relationships to better understand their emotional dynamics.} \\
\texttt{3. Be more nuanced in my emotional analysis, avoiding oversimplification and taking into account the complexity of human emotions.} \\
\\ 
\texttt{By following this plan, I aim to improve my ability to predict the likely emotional responses of characters in dialogue, and to provide more accurate and insightful analysis.} \\
            }
        }
    }
    & \\ \\

\end{supertabular}
}

\end{document}
